\documentclass{beamer}

\usepackage{beamerthemeSingapore}
\usepackage{graphicx}
%\usepackage{here}
\usepackage{url}
\usepackage[brazil]{babel}
\usepackage[utf8]{inputenc}
\usepackage[tight,raggedright]{subfigure}
\usepackage[alf]{abntex2cite}
\usepackage{listings}
\usepackage{minted}
\usepackage{amsmath}
\usepackage{pdfpages}
% Please add the following required packages to your document preamble:
\usepackage{booktabs}

% TODO TODO TODO TODO TODO TODO TODO
\title{03 - Recursão}
\author{Prof. Alex Kutzke}
\date{Maio de 2021}

\begin{document}

	% ==========================================
	\begin{frame}
		\begin{center}
			% TODO TODO TODO TODO TODO TODO
			\large{Setor de Educação Profissional e Tecnológica}\\
			Tecnologia em Análise e Desenvolvimento de Sistemas\\
			\vspace{1cm}			
			\small{DS141 - Estruturas de Dados II}
		\end{center}
		\titlepage
	\end{frame}
	% ==========================================

	% ==========================================
	\begin{frame}[c]{Roteiro}
		\tableofcontents
	\end{frame}
	% ==========================================

	\section{Definição}

  \begin{frame}[c]{Recursão}
  	\begin{itemize}
  		\item Fundamental na matemática e na computação;

			\item Definição: ``Algo definido em \textbf{termos de si próprio}'';

			\begin{itemize}
				\item Matemática: Uma \textbf{função recursiva} é aquela que é definida em termos de si própria;
				\item Computação: Um \textbf{programa recursivo} é aquele que ``chama'' (executa) a si mesmo.
			\end{itemize}

			\item \textbf{Algoritmo recursivo:} aquele que resolve um problema a partir da solução de uma ou mais instâncias \textit{menores} do mesmo problema \cite{sedgewick2014algorithms};
			\begin{itemize}
				\item Em $C$, utilizamos funções recursivas para implementar programas recursivos.
			\end{itemize}
  	\end{itemize}
  \end{frame}

	\begin{frame}[c]{Relação de Recorrência}
		\begin{itemize}
			\item Relações de recorrência são funções definidas recursivamente (matemática);

			\item Uma relação de recorrência é definida por duas partes:
			\begin{itemize}
				\item conjunto de valores iniciais, ou bases; e
				\item valores determinados em termos da própria função para valores \textbf{menores} de ``entrada''.
			\end{itemize}

			\item Talvez a mais conhecida seja a função \textit{fatorial}:
		\end{itemize}

		\begin{center}
			\[ N! =
			\left \{
			\begin{array}{rc}
				1,&\mbox{se}\quad N = 0,\\
				N * (N-1)!,&\mbox{se}\quad N \ge 1.
			\end{array}
			\right .
			\]
		\end{center}
	\end{frame}

\lstset{language=C,
								 basicstyle=\ttfamily,
                 keywordstyle=\color{blue}\ttfamily,
                 stringstyle=\color{red}\ttfamily,
                 commentstyle=\color{green}\ttfamily,
                 morecomment=[l][\color{magenta}]{\#}}

	\begin{frame}[fragile,c]{Implementação Recursiva Fatorial}
	
		\begin{minted}[
				frame=lines,
				framesep=2mm,
				baselinestretch=1.2,
				%fontsize=\footnotesize,
				linenos
			]{c}
int fat(int n)
{
	if(n == 1) return(1);
	return(n * fat(n-1));
}
		\end{minted}

		\begin{itemize}
			\item É sempre possível transformar um algoritmo recursivo em um não-recursivo:
		\end{itemize}
	
		\begin{minted}[
				frame=lines,
				framesep=2mm,
				baselinestretch=1.2,
				%fontsize=\footnotesize,
				linenos
			]{c}
for(t=1,i=1;i<=n;i++) t*=i;
		\end{minted}
		\begin{itemize}
			\item O inverso também é verdadeiro: é possível transformar qualquer algotirmo não-recursivo que envolva laços em algoritmos recursivos.
		\end{itemize}
	\end{frame}

	\begin{frame}[c]{Análise Algoritmo Fatorial}
		\begin{itemize}
			\item O algoritmo de fatorial apresentado ilustra as principais propriedades da recursão:
			\begin{itemize}
				\item Possui uma condição de parada (término);
				\item Executa ele mesmo para valores cada vez menores.
			\end{itemize}

			\item Indução matemática pode ser utilizada para provar que o programa funciona como esperado:

			\begin{itemize}
				\item Computa 0! (base da indução);
				\item Sobre a suposição de que computa $k!$, para $k < N$ (hipótese indutiva), temos que o algoritmo computa $N!$.
			\end{itemize}
			\item O mesmo raciocínio pode ser utilizado para resolver problemas mais complexos.
		\end{itemize}
	\end{frame}

	\section{Propriedades}

	\begin{frame}[c]{Condições}
		Em resumo, uma função recursiva precisa:
		\begin{itemize}
			\item Resolver pelo menos um caso base;
			\item Cada chamada recursiva deve envolver valores menores de argumentos.
		\end{itemize}

		Esses itens permitem acreditar que uma prova por indução é possível. Além de serem ótimos guias de implementação para algoritmos recursivos.
	\end{frame}

	\begin{frame}[c]{Vantagens e desvantagens}
		\begin{itemize}
			\item Vantagens:
			\begin{itemize}
				\item Deixa o código mais simples e mais compacto;
				\begin{itemize}
					\item Nem sempre significa que é mais fácil de entender, mas sim que segue a definição matemática.
				\end{itemize}
				\item Fácil implementação (seguindo a recorrência matemática).
			\end{itemize}
			\item Desvantagens:
			\begin{itemize}
				\item Podem se tornar impraticáveis em termos de custo computacional;
				\begin{itemize}
					\item É necessário fugir de implementações ruins.
				\end{itemize}
				\item Seu desempenho depende da implementação de pilhas do sistema operacional e da \textit{profundidade da recursão};
				\item Em geral utilizam mais memória e são de difícil depuração.
			\end{itemize}
		\end{itemize}
	\end{frame}
	
	\section{Exemplos}

	\begin{frame}[fragile, c]{Qual o problema dessa função?}
		\begin{minted}[frame=lines,framesep=2mm,baselinestretch=1.2,linenos]{c}
int puzzle(int n)
{
	if(n == 1) return(1);
	if(n % 2 == 0)
		return(puzzle(n/2));
	else return(puzzle(3*n+1));
}
		\end{minted}
	\end{frame}

	\begin{frame}[c]{Máximo Divisor Comum}
		\begin{itemize}
			\item Algoritmo de euclides:
			\begin{itemize}
				\item Baseia-se na seguinte observação:
				\begin{center}
					O máximo divisor comum de dois números inteiros $x$ e $y$, com $x > y$, é o mesmo que o máximo divisor comum de $y$ e $x \% y$.

					\vspace{10pt}
					\pause
					Um número $t$ divide ambos $x$ e $y$ se, e somente se, $t$ divide ambos $y$ e $x \% y$, pois $x$ é igual a $x \% y$ mais um múltiplo de $y$.
				\end{center}
			\end{itemize}
		\end{itemize}
	\end{frame}

	\begin{frame}[fragile,c]{Algoritmo de Euclides Recursivo}
		\begin{minted}[frame=lines,framesep=2mm,baselinestretch=1.2,linenos]{c}
int mdc(int m, int n)
{
	if(n == 0) return(m);
	return(mdc(n, m % n));
}
		\end{minted}
	\end{frame}

	\begin{frame}[c]{Avaliação de Expressões Pré-fixas}
		Como seria um algoritmo recursivo para expressões do tipo:\\

		$*$ $+$ $7$ $*$ $*$ $4$ $6$ $+$ $8$ $9$ $5$ ?
	\end{frame}

	\begin{frame}[fragile,c]{Avaliador Pré-fixo}
		\begin{minted}[frame=lines,framesep=2mm,baselinestretch=1.2,linenos]{c}
char a[100]; int i;
int eval(){
	int x = 0;
	while(a[i] == ' ') i++;
	if(a[i] == '+') { i++;
		return(eval() + eval()); }
	if(a[i] == '*') { i++;
		return(eval() * eval()); }
	while((a[i] >= '0') && (a[i] <= '9'))
		x = 10*x + (a[i++]-'0');
	return(x);
}
		\end{minted}
	\end{frame}

	\begin{frame}[c]{Profundidade da Recursão}
		\begin{itemize}
			\item A chamada de uma função carrega consigo um custo computacional ``escondido'':
			\begin{itemize}
				\item Pilha de chamadas de função.
			\end{itemize}
			\item A \textbf{profundidade} de uma recursão é a quantidade de chamadas recursivas que a função realiza.
		\end{itemize}
	\end{frame}

	\begin{frame}[fragile,c]{Algoritmo Fibonacci Recursivo}
		\begin{minted}[frame=lines,framesep=2mm,baselinestretch=1.2,linenos]{c}
int fib(int n)
{
	if(n == 0) return(0);
	if(n == 1) return(1);
	return(fib(n-1) + fib(n-2));
}
		\end{minted}
		\begin{center}
			Quantas chamandas recursivas esse algoritmo faz quando $n = 25$?
		\end{center}

		\begin{center}
			O que se pode dizer sobre a complexidade desse algoritmo?
		\end{center}
	\end{frame}

	\section{Recursão e TADs}

	\begin{frame}[c]{Recursão e TADs}
		\begin{itemize}
			\item TADs baseadas em nós e ponteiros possuem, em geral, um caráter recursivo;

			\item Como contar o número de elementos de uma lista utilizando uma função recursiva?

			\item Como encontrar o maior elementa da lista?
		\end{itemize}
	\end{frame}

	\begin{frame}[fragile,c]{O que esse algoritmo faz?}
		\begin{minted}[fontsize=\scriptsize,frame=lines,framesep=2mm,baselinestretch=1.2,linenos]{c}
int aaaa(int a[],int l, int r){
	int u,v;
	if(l == r) return(a[l]);
	int m = (l+r)/2;
	u = aaaa(a,l,m);
	v = aaaa(a,m+1,r);
	if(u > v) return(u); else return(v);
}
...
aaaa(v,0,tam_v);
		\end{minted}

		\pause

		\begin{itemize}
			\item Encontrar o máximo do vetor;
			\item Abordagem dividir para conquistar.
		\end{itemize}

		\begin{center}
			Uma função recursiva que divide um problema de tamanho $N$ em duas partes independentes (não vazias) que são resolvidas recursivamente, faz menos de $N$ chamadas recursivas.
		\end{center}
	\end{frame}

	\begin{frame}[fragile,c]{Programação Dinâmica}
		\begin{itemize}
			\item Técnica de programação;
			\item Armazenar valores já calculados para chamadas recursivas;
		\end{itemize}
		\begin{minted}[fontsize=\scriptsize,frame=lines,framesep=2mm,baselinestretch=1.2,linenos]{c}
int knownF[100];
int fib(int n)
{
	int t;

	if(knownF[n] != 0) return knownF[n];
	if(n == 0) return(0);
	if(n == 1) return(1);
	t = fib(n-1) + fib(n-2);
	return knownF[n] = t;
}
		\end{minted}
		\begin{itemize}
			\item Quantas chamadas recursivas esse algoritmo faz?
		\end{itemize}
	\end{frame}
	%--- Next Frame ---%

	% ==========================================
	% \begin{frame}
	% 	\frametitle{Título}
	% 	\begin{itemize}
	% 		\item Item bla bla bla
	% 	\end{itemize}
	%	  \begin{block}{Título de bloco}
	%		 	Conteúdo do bloco
	%   \end{block}
	% \end{frame}
	% ==========================================

	% ==========================================
	% Exemplo de frame com figura
	% \begin{frame}
	% 	\frametitle{Título}
	%	  \begin{figure}[h]
	%	  	\centering
	%	  	\includegraphics[width=10cm]{paht.png}
	%			\caption{Bla bla bla ...}
	%	  	\label{label}
	%	  \end{figure}
	% \end{frame}
	% ==========================================

	% ==========================================
	% Exemplo de frame com algoritmo
	% \begin{frame}[fragile]
	% 	\frametitle{Título}
	% 	\begin{lstlisting}[language=C]
	% 		#include <stdio.h>
	%
	% 		int main()
	% 		{
	% 		  int a[5];
	%
	% 		  a[0] = 30;
	% 		  a[1] = 25;
	% 		  a[2] = 5;
	% 		  a[3] = 100;
	% 		  a[4] = 1;
	%
	% 		  printf("%d\n",a[1]);
	% 		}
	% 	\end{lstlisting}
	% \end{frame}
	% ==========================================

	% ==========================================
	% Frame para referências
	\begin{frame}[allowframebreaks]
		\frametitle{Referências utilizadas na apresentação}
		\bibliography{referencias}
	\end{frame}
  % ==========================================

\end{document}
